% Generated draft: requires mathematical and source review.
\documentclass[11pt]{article}
\usepackage[margin=25mm]{geometry}
\usepackage{fontspec}
\setmainfont{Latin Modern Roman}
\usepackage{amsmath,amssymb,mathtools,unicode-math}
\setmathfont{Latin Modern Math}
\usepackage{enumitem,needspace,xurl,hyperref}
\setlength{\parindent}{0pt}
\setlength{\parskip}{5pt}
\setlength{\emergencystretch}{3em}
\newcommand{\N}{\mathbb{N}}
\newcommand{\Z}{\mathbb{Z}}
\newcommand{\Q}{\mathbb{Q}}
\newcommand{\R}{\mathbb{R}}
\newcommand{\C}{\mathbb{C}}
\newcommand{\F}{\mathbb{F}}
\newcommand{\A}{\mathbb{A}}
\newcommand{\PP}{\mathbb{P}}
\newcommand{\E}{\mathbb{E}}
\newcommand{\eps}{\varepsilon}
\newcommand{\dd}{\,\mathrm{d}}
\begin{document}
\section*{UnsolvedMath ID: 11100012}
\textbf{Source problem code:} catalogue entry
\subsection{Definitions and mathematical statement}
Fix a prime $p$ and an intended chromatic height $n\geq1$. In the customary interpretation, the topological Picard group $\operatorname{Pic}_n$ consists of equivalence classes of invertible $K(n)$-local spectra: $X$ is invertible if there is $Y$ with $L_{K(n)}(X\wedge Y)\simeq L_{K(n)}\mathbb S$. The group operation is the localized smash product.

Hovey's Morava K- and E-theory item 2 asks to show that the Picard group is finitely generated ``over the $p$-adics.'' Here $\Z_p=\varprojlim_r\Z/p^r\Z$. For a specified $\Z_p$-module $P$, the precise finite-generation assertion would be
\[
 \exists r<\infty\quad\exists x_1,\ldots,x_r\in P,
 \qquad P=\Z_p x_1+\cdots+\Z_p x_r.
\]
If $P$ is a compact module this can be stated as a continuous surjection $\Z_p^r\twoheadrightarrow P$.

The source does not define the Picard group or its $\Z_p$-module structure in this excerpt. This matters: a full Picard group may have prime-to-$p$ finite torsion, which cannot occur in a $\Z_p$-module. The assertion may therefore concern a pro-$p$ component, a finite-index $p$-adic part, or finite generation with additional finite torsion allowed. These readings must be resolved, not silently identified.

The source also asks for algebraic understanding through profinite group cohomology and says it did not believe finite generation known even for the algebraic Picard group. In the usual Morava-$E$ setting that algebraic object concerns invertible graded $E_{n,*}$-modules with continuous semilinear stabilizer action, with a degree-zero part described by continuous first cohomology of units; the grading and stabilizer conventions are not supplied. The registry reports low-height algebraic computations (BSSW, 2024), while saying arbitrary-height topological finite generation is not fully settled. The historical uncertainty and this later report are distinct.
\subsection{Short English statement}
Show suitable $p$-adic finite generation of the chromatic Picard group, and investigate its algebraic group-cohomology counterpart. The full group, finite torsion and module conventions must be specified.
\subsection{Sources}
\begin{itemize}
\item Local registry: \url{data/unsolvedmath/problems.json}, permanent UnsolvedMath ID $11100012$, exact source code \texttt{catalogue entry}, attributed to Mark Hovey.
\item \emph{Hovey's Algebraic Topology Problem List}, surviving mirror, Morava K- and E-theory item 2: \url{https://www-users.cse.umn.edu/~tlawson/hovey/}. Supplied access date: 2026-07-29; extraction: html-ordered-list; 65 top-level items in nine category pages.
\item The registry cites ``BSSW 2024'' for low-height algebraic Picard calculations but supplies no full author names, title, identifier, heights or prime restrictions. No other precise primary citation is included.
\end{itemize}
\end{document}
